\section{CodeClaimVerifier Design}
\label{sec:design}

\subsection{Architecture Overview}

CCV operates in three stages (Figure~\ref{fig:architecture}):

\begin{enumerate}
    \item \textbf{Typed Claim Extraction} (1 LLM call): A structured-output LLM call decomposes reasoning into a list of typed claims, each with a claim type, parameters, and source sentence.
    \item \textbf{Deterministic Verification} (0 LLM calls): Each claim dispatches to a type-specific verifier that checks the assertion against the codebase using grep, file reads, or lockfile parsing.
    \item \textbf{Confidence Calibration}: Verification results are aggregated into a calibrated confidence score with an actionable verdict.
\end{enumerate}

\begin{figure}[t]
\centering
\small
\begin{verbatim}
Agent Reasoning (text) + Evidence (dict)
    |
    v
[Phase 1: Typed Claim Extraction]  -- 1 LLM call
    |
    v
List[TypedClaim]  -- type, params, source
    |
    v
[Dependency Graph + Topological Sort]
    |
    v
[Phase 2: Deterministic Verification]
    |   grep, read, parse per claim type
    |   language-aware patterns
    |   path traversal prevention
    v
List[VerifiedClaim]  -- verdict, confidence
    |
    v
[SUSPECT Propagation]
    |
    v
[Confidence Calibration]  -- pure math
    |
    v
VerificationReport  -- rate, action, per_claim
\end{verbatim}
\caption{CCV architecture. One LLM call for extraction, zero for verification.}
\label{fig:architecture}
\end{figure}

\subsection{Claim Taxonomy}

CCV defines 14 claim types organized into four categories (Table~\ref{tab:taxonomy}). Each type has a defined parameter schema, verification method, and method confidence level.

\begin{table*}[t]
\centering
\caption{CCV claim taxonomy. Confidence values are engineering estimates; empirical calibration is evaluated in Section~\ref{sec:evaluation}.}
\label{tab:taxonomy}
\small
\begin{tabular}{llllr}
\toprule
\textbf{Category} & \textbf{Claim Type} & \textbf{Parameters} & \textbf{Verification Method} & \textbf{Conf.} \\
\midrule
\multirow{4}{*}{File/Path}
& \textsc{File\_Exists} & path & \texttt{os.path.isfile()} & 0.99 \\
& \textsc{Line\_Content} & path, line, expected & File read + string match & 0.95 \\
& \textsc{File\_Classification} & path, category & Path pattern regex & 0.80 \\
& \textsc{Generated\_Or\_Vendored} & path, expected & Header markers + vendor dirs & 0.85 \\
\midrule
\multirow{3}{*}{Function}
& \textsc{Function\_Exists} & name, file (opt.) & Language-aware def grep & 0.85 \\
& \textsc{Function\_Called} & name, expected & Call-site grep minus defs & 0.65 \\
& \textsc{Has\_Callers} & name, expected & Call-site grep minus defs & 0.65 \\
\midrule
\multirow{4}{*}{Dependency}
& \textsc{Import\_Exists} & module, file (opt.) & Language-aware import grep & 0.80 \\
& \textsc{Package\_Version} & package, version & Lockfile parse & 0.90 \\
& \textsc{Dependency\_Type} & package, type & Manifest parse & 0.85 \\
& \textsc{Cve\_Affects\_Version} & cve, package, version & Advisory DB (unimpl.) & --- \\
\midrule
\multirow{3}{*}{Code}
& \textsc{Absence} & pattern, scope & Scoped grep (negated) & 0.60 \\
& \textsc{Mitigation\_Exists} & file, line & File read + non-empty check & 0.70 \\
& \textsc{Entry\_Point} & type, location & Framework pattern grep & 0.65 \\
\bottomrule
\end{tabular}
\end{table*}

\textbf{Language awareness.} Function definition and import patterns are language-specific. CCV supports Python, Go, TypeScript/JavaScript, Java, C/C++, and Rust. Language is detected from the file extension of the finding under analysis, not the repository's dominant language.

\textbf{Absence claims.} Unlike Claimify, which filters out negation claims, CCV treats absence as a first-class claim type. The \textsc{Absence} verifier uses scoped grep: if the pattern is not found in the specified scope (file, directory, or repo), the claim is \textsc{Verified}. This captures assertions like ``no input sanitization exists'' or ``the function has no callers.'' The confidence is lower (0.60) because grep cannot detect aliases, dynamic dispatch, or language-level indirection.

\subsection{Claim Chaining}

LLM claims often have implicit dependencies. A \textsc{Line\_Content} claim that ``line 42 contains \texttt{torch.load()}'' implicitly depends on the file existing. If the file doesn't exist, the line content claim is meaningless regardless of whether the verifier would match the pattern.

CCV infers these dependencies from shared parameters using seven built-in rules (Table~\ref{tab:chaining}). When a dependent claim exists but its prerequisite was not explicitly extracted, CCV synthesizes the prerequisite and verifies it.

\begin{table}[t]
\centering
\caption{Dependency inference rules. Source and target parameters define how synthesized claims are constructed.}
\label{tab:chaining}
\small
\begin{tabular}{lll}
\toprule
\textbf{Dependent} & \textbf{Prerequisite} & \textbf{Params} \\
\midrule
\textsc{Line\_Content} & \textsc{File\_Exists} & path$\to$path \\
\textsc{Generated\_Or\_Vendored} & \textsc{File\_Exists} & path$\to$path \\
\textsc{Function\_Exists} & \textsc{File\_Exists} & file$\to$path \\
\textsc{Function\_Called} & \textsc{Function\_Exists} & name$\to$name \\
\textsc{Has\_Callers} & \textsc{Function\_Exists} & name$\to$name \\
\textsc{Import\_Exists} & \textsc{File\_Exists} & file$\to$path \\
\textsc{Mitigation\_Exists} & \textsc{File\_Exists} & file$\to$path \\
\bottomrule
\end{tabular}
\end{table}

\textbf{Resolution algorithm.} Claims are verified in topological order. If all prerequisites of a given type are \textsc{Refuted}, the dependent is marked \textsc{Suspect} with reduced confidence contribution. CCV uses ANY-match semantics: if at least one prerequisite of a type is \textsc{Verified}, the dependent is not flagged.

\textbf{Synthesized claims} are marked as such and excluded from report metrics to prevent phantom claims from distorting confidence scores.

\subsection{Confidence Calibration}

CCV computes a weighted verification rate:

\begin{equation}
\text{rate} = \frac{\sum_{c \in V^+} w(c) \cdot \alpha_c}{\sum_{c \in V} \alpha_c}
\label{eq:rate}
\end{equation}

where $V$ is the set of verifiable (non-\textsc{Unverifiable}) claims, $V^+$ is the set of \textsc{Verified} claims, $\alpha_c$ is the method confidence, and $w(c) = 0.5$ if $c$ is \textsc{Suspect}, $w(c) = 1.0$ otherwise. This asymmetric weighting ensures that \textsc{Suspect}-\textsc{Verified} claims lower the rate (they contribute full confidence to the denominator but half to the numerator).

The rate maps to actions:

\begin{itemize}
    \item $\text{rate} \geq 0.8$: \textsc{Boost} -- increase confidence in the LLM's conclusion
    \item $0.5 \leq \text{rate} < 0.8$: \textsc{Flag} -- flag for human review
    \item $\text{rate} < 0.5$: \textsc{Override} -- override the LLM's conclusion
\end{itemize}

\subsection{Extensibility}

CCV supports user-defined claim types via a registration API. Users provide a verifier function, extraction hint (injected into the LLM prompt), and optional dependency rules. Custom types receive the same caching, chaining, and calibration as built-in types.

\subsection{Caching}

For batch verification (multiple findings against the same codebase), CCV uses two cache layers:

\begin{itemize}
    \item \textbf{Grep cache}: Thread-safe via \texttt{contextvars}, keyed by (pattern, path, fixed). Returns defensive copies.
    \item \textbf{Verifier cache}: Keyed by (type, frozen\_params, repo, language). Stores bare results (pre-chaining) and returns copies via \texttt{dataclasses.replace()} for independent SUSPECT marking.
\end{itemize}

Both caches are per-call (no staleness across invocations) and support concurrent use.
